\chapter{Monte Carlo、bootstrap 与方差缩减}
\label{ch:monte-carlo}

\begin{introduction}
  \item 先修：第 9 章的大数定律与中心极限定理，第 14 章的浮点误差和容差。
  \item 学习产出：能从期望或积分推导 Monte Carlo 估计量及误差预算，使用独立 oracle 验收模拟，解释并实现控制变量，构造 bootstrap 抽样分布，并判断 IID bootstrap 何时失效。
  \item 研究连接：衍生品定价、风险度量、策略不确定性与执行仿真都需要把抽样误差、模型误差、数值误差和实现错误分开记账。
\end{introduction}

\section{模拟解决的不是“不会积分”，而是高维期望}

很多量化任务都可写成某个分布下的期望
\[
  \mu=\mathbb E_P[g(X)]=\int g(x)\,dP(x).
\]
例如路径依赖期权的贴现收益、组合损失超过阈值的概率、复杂执行策略的成本，往往没有方便的闭式解。Monte Carlo 用可抽样的 $P$ 把积分变成样本平均；bootstrap 则用观测得到的经验分布近似未知的抽样分布。两者都“重复抽样”，但回答的问题不同：前者近似给定模型下的期望，后者近似统计量在重复采样下的不确定性。

在运行随机数之前，先写清五件事：目标分布 $P$、被积函数或统计量、样本之间的依赖结构、独立 oracle、允许的误差。若这些对象没有定义，增加模拟次数只会更稳定地计算错误问题。

\begin{example}
设 $U\sim\operatorname{Unif}(0,1)$，目标为 $\mu=\mathbb E[U^2]$。解析积分立即给出 $\mu=1/3$；这不是 Monte Carlo 程序的一部分，而是验收程序的独立 oracle。真实项目没有闭式解时，也应先寻找可退化到闭式解的参数、小规模可枚举状态或第二种独立算法。
\end{example}

\section{Monte Carlo 估计量与 $N^{-1/2}$ 误差率}

设 $X_1,\ldots,X_N$ 独立同分布，$Y_i=g(X_i)$，并假设
$\mu=\mathbb E[Y_i]$、$0<\sigma^2=\operatorname{Var}(Y_i)<\infty$。定义
\[
  \hat\mu_N=\frac1N\sum_{i=1}^N Y_i.
\]
线性期望给出无偏性：
\[
  \mathbb E[\hat\mu_N]
  =\frac1N\sum_{i=1}^N\mathbb E[Y_i]=\mu.
\]
方差推导必须显式使用独立性：
\[
\begin{aligned}
  \operatorname{Var}(\hat\mu_N)
  &=\frac1{N^2}\sum_{i=1}^N\sum_{j=1}^N
    \operatorname{Cov}(Y_i,Y_j)\\
  &=\frac1{N^2}\sum_{i=1}^N\sigma^2
   =\frac{\sigma^2}{N}.
\end{aligned}
\]
因此理论标准误为
\[
  \operatorname{SE}(\hat\mu_N)=\frac{\sigma}{\sqrt N}.
\]
把标准误减半通常需要四倍样本；固定 seed 只能复现实验，不能改善误差率。若样本相关，上式中的非对角协方差不再消失，有效样本量会小于名义的 $N$。

\subsection{从误差目标反推样本量}

在中心极限定理近似可信且使用双侧正态区间时，希望半宽不超过 $\varepsilon$，可先用
\[
  N\gtrsim \left(\frac{z_{1-\alpha/2}\sigma}{\varepsilon}\right)^2
\]
做预算。对 $U^2$，解析可得
\[
  \sigma^2=\mathbb E[U^4]-\mathbb E[U^2]^2
  =\frac15-\frac19=\frac4{45}.
\]
若希望约 95\% 区间的半宽为 $10^{-3}$，代入 $z_{0.975}\approx1.96$ 得到约
$3.42\times10^5$ 个独立样本。这一数量级提醒我们：Monte Carlo 的通用性以较慢的平方根收敛为代价。

\subsection{统计容差不是固定小数位}

证据实验使用固定随机源，却不锁死某次估计的小数。它要求
\[
  |\hat\mu_N-1/3|\le 4\sqrt{(4/45)/N},
\]
并用 $N=256$ 与 $N=4096$ 的多次独立重复检查 RMSE 比率接近
$\sqrt{4096/256}=4$。四个标准误不是“永远正确”的保证，而是在有限方差与近似正态误差下声明的显式验收预算。重尾、稀有事件、强自相关和自适应停止都要求另外推导。

\section{建立完整误差预算}

模拟输出与真值的差异不只来自 Monte Carlo 方差。一份可审计报告应至少区分：
\begin{enumerate}[nosep,leftmargin=*]
  \item \textbf{模型误差：}目标分布、相关结构或市场机制本身不正确；
  \item \textbf{离散化偏差：}时间步长、网格或近似动态改变了目标；
  \item \textbf{抽样误差：}有限随机样本导致的方差；
  \item \textbf{数值误差：}浮点归约、病态线性系统或权重上溢；
  \item \textbf{实现错误：}采样器、支付函数、单位或折现代码与定义不一致。
\end{enumerate}
解析解或可枚举小问题先排除实现错误；步长外推检查离散化偏差；独立 seeds 或不重叠子流评估抽样不确定性；第 14 章的高精度与尺度实验约束数值误差。只展示一条“看起来很平滑”的收敛曲线会把随机波动与选择性汇报藏起来。

对置信区间还应说明使用的是已知方差、样本方差、渐近正态近似还是非渐近界。估计对象若是分位数、最大回撤或罕见违约概率，均值中心极限定理不能自动证明同一种区间有效。

\section{伪随机数、seed 与并行子流}

伪随机数生成器从有限内部状态确定性地产生序列。seed 是复现协议的一部分，但不是质量证明。研究报告应记录生成器、seed 或 seed 序列、样本数量、并行划分及软件版本；测试与最终研究结果宜使用不同但预先登记的随机流，避免反复换 seed 直到结论好看。

并行模拟不能让多个进程从同一内部状态开始，否则会重复路径并夸大有效样本量。更可靠的做法是使用库提供的子流或 seed sequence 确定性地产生不重叠状态，并固定“任务编号到子流”的映射。并行归约还会改变浮点加法顺序；若差异影响结论，应使用稳定求和或把归约顺序纳入误差预算。

\begin{note}
复现不是把所有研究永久绑定到一个 seed。更好的证据是：一个登记 seed 用于精确复跑，多个独立 seeds 用于评估结论对随机源的敏感性，解析或枚举 oracle 用于判断程序是否计算了正确对象。
\end{note}

\section{控制变量：用已知期望换取更小方差}

若 $h(X)$ 的期望 $m=\mathbb E[h(X)]$ 已知，可构造
\[
  Y_\beta=g(X)-\beta\{h(X)-m\}.
\]
只要 $m$ 正确，就有 $\mathbb E[Y_\beta]=\mathbb E[g(X)]$。展开方差：
\[
  \operatorname{Var}(Y_\beta)
  =\operatorname{Var}(g)-2\beta\operatorname{Cov}(g,h)
   +\beta^2\operatorname{Var}(h).
\]
这是关于 $\beta$ 的凸二次函数。令导数为零得到
\[
  \beta^*=\frac{\operatorname{Cov}(g,h)}{\operatorname{Var}(h)}.
\]
代回可得最小方差
\[
  \operatorname{Var}(g)(1-\rho_{g,h}^2).
\]
控制变量越相关，潜在收益越大；但相关不代表已知期望，错误的 $m$ 会直接引入偏差。

\subsection{完整例子：$U^2$ 用 $U$ 作控制变量}

取 $g(U)=U^2$、$h(U)=U$，有
\[
  m=\frac12,\quad
  \operatorname{Cov}(U^2,U)=\frac14-\frac13\frac12=\frac1{12},\quad
  \operatorname{Var}(U)=\frac1{12}.
\]
故解析最优系数为 $1$。调整量为 $U^2-U+1/2$，其方差
\[
  \operatorname{Var}(U^2-U)=\frac1{180},
\]
而原方差为 $4/45$，理论方差缩减倍数为 $16$。证据程序同时检查两件事：调整后均值仍在 $1/3$ 的解析标准误预算内，样本方差缩减超过保守下界 10。若把已知均值从 $1/2$ 错写为 $0.6$，方差完全不变却引入 $0.1$ 的偏差；专门的负例必须拒绝这种“低方差错误答案”。

实际问题常需从数据估计 $\beta^*$。同一批样本既估系数又汇报收益可能产生小样本偏差或乐观评价；可使用独立 pilot 样本、样本分割或交叉拟合，并把估计成本计入总预算。

\section{其他方差缩减方法及其边界}

\subsection{反变量与共同随机数}

对 $U\sim\operatorname{Unif}(0,1)$，$1-U$ 仍服从同一分布。把
$g(U)$ 与 $g(1-U)$ 成对平均，在单调 $g$ 下常产生负相关并降低方差。比较两个策略时，共同随机数让两者经历相同市场路径，目标是降低“差值估计”的方差，而不是分别让每个估计更精确。两种方法都需要检查所构造边缘分布没有改变。

\subsection{分层抽样}

把样本空间分成互斥层 $A_k$，则
\[
  \mu=\sum_k P(A_k)\,\mathbb E[g(X)\mid A_k].
\]
在每层独立抽样并按真实层概率加权，可避免普通随机抽样遗漏重要区域。若层权重使用了样本频率而不是真实概率，估计对象会改变；最优分配还依赖层内方差与抽样成本。

\subsection{重要性抽样}

若从提议分布 $Q$ 抽样且 $P$ 对 $Q$ 绝对连续，则
\[
  \mathbb E_P[g(X)]
  =\mathbb E_Q\!\left[g(X)\frac{dP}{dQ}(X)\right].
\]
它适合把更多样本放到稀有但重要的区域。失败边界也很明确：$Q$ 没有覆盖 $P$ 的重要支撑会产生偏差；权重重尾会让方差无穷或由少量路径支配。报告应包含权重分布、有效样本量和提议分布敏感性，而不只给更小的某次误差。

\section{Bootstrap：条件于样本的重抽样}

给定观察 $X_1,\ldots,X_n$，经验分布为
\[
  \hat F_n=\frac1n\sum_{i=1}^n\delta_{X_i}.
\]
IID bootstrap 从 $\hat F_n$ 有放回抽取大小为 $n$ 的样本
$X_1^*,\ldots,X_n^*$，计算目标统计量 $T_n^*=T(X_1^*,\ldots,X_n^*)$，再用条件分布
\[
  \mathcal L(T_n^*\mid X_1,\ldots,X_n)
\]
近似原统计量的抽样分布。它不是从真实总体重新收集数据，而是 plug-in 原理：用 $\hat F_n$ 代替未知的 $F$。

\subsection{可枚举 oracle：样本 $(1,2,4)$}

大小为 3 的有序重抽样共有 $3^3=27$ 个。若用 $S^*$ 表示重抽样和，其可能取值及出现次数为
\[
\begin{array}{c|rrrrrrrrr}
S^*&3&4&5&6&7&8&9&10&12\\\hline
\text{次数}&1&3&3&4&6&3&3&3&1
\end{array}
\]
重抽样均值是 $\bar X^*=S^*/3$，所以
\[
  \mathbb E_*(\bar X^*)=\frac73,
  \qquad
  \operatorname{Var}_*(\bar X^*)
  =\frac1{27}\sum_s n_s\left(\frac{s}{3}-\frac73\right)^2
  =\frac{14}{27}.
\]
这里星号表示“条件于原样本的 bootstrap 世界”。穷举 27 个重抽样均值和 bootstrap 方差为 $14/27$ 构成精确 oracle；随机重复只应逼近它。

\subsection{从重抽样分布构造不确定性}

常见输出包括 bootstrap 标准误、偏差估计和置信区间。percentile 区间直接取 $T_n^*$ 的分位数；basic 区间围绕原估计反射 bootstrap 分位数；normal 区间依赖近似对称与标准误；BCa 还修正偏差和加速度。它们不是可随意互换的绘图选项：统计量是否平滑、样本量、偏斜与边界会改变覆盖率。

模拟次数 $B$ 也有自己的 Monte Carlo 误差。若关注 2.5\% 分位数，尾部有效重抽样远少于 $B$；应检查分位数随 $B$ 和 seed 的稳定性。把 $B$ 从一千增至一百万不会修复错误的 IID 假设。

\section{Bootstrap 何时失效}

IID bootstrap 的关键前提是观察近似可交换，且目标统计量对经验分布足够规则。以下情形需要额外方法或更弱结论：
\begin{itemize}[nosep,leftmargin=*]
  \item \textbf{时间依赖：}逐日重抽样破坏自相关、波动聚集和持仓路径；可考虑移动块、圆形块或平稳 bootstrap，并解释 block 长度。
  \item \textbf{横截面簇：}同一公司、行业或日期内误差相关；应按独立簇重抽样，而非逐行重抽样。
  \item \textbf{极值与边界：}样本最大值无法产生超出已观测最大值的新尾部，某些极值统计量的普通 bootstrap 不一致。
  \item \textbf{重尾与小样本：}经验分布可能根本没观察到决定风险的尾部区域；区间看似精确但遗漏总体不确定性。
  \item \textbf{模型选择：}先在全样本挑因子或调参，再只 bootstrap 最终模型，会忽略选择步骤；重抽样流程必须包含真实研究管线。
\end{itemize}
对时间序列收益直接做 IID bootstrap 后得到很窄的 Sharpe 区间，并不证明策略稳定。应先定义独立单位和目标统计量，再比较依赖感知方案、不同 block 规则、时间外样本与模拟覆盖率。证据程序在数据声明为有依赖时稳定拒绝 IID bootstrap，不悄悄输出一个错误区间。

\section{Python 证据程序如何验收结论}

权威 Jupytext 源 \path{notebooks/upper/ch15_monte_carlo.py} 依次执行四个层次：
\begin{enumerate}[nosep,leftmargin=*]
  \item 用解析 $1/3$ 与 $4/45$ 验收样本均值及标准误；
  \item 用两个样本量和一千次独立重复验收平方根误差率；
  \item 用解析 $\beta^*=1$、方差 $1/180$ 验收控制变量的均值与方差；
  \item 用有理数穷举 27 个 bootstrap 状态，不从模拟程序生成期望答案。
\end{enumerate}
正常输出只暴露稳定类别，例如 \texttt{error<=4se} 与 \texttt{sample\_vrf>=10}，不把平台相关的随机小数写成契约。两个负例分别把控制变量中心改错、把有依赖数据声明交给 IID bootstrap；它们必须给出稳定诊断。这样固定 seed、理论容差、独立 oracle 和失败协议各自承担不同职责。

\section{研究报告与面试检查表}

面对“你如何验证 Monte Carlo 或 bootstrap”时，可以按以下顺序回答：
\begin{enumerate}[nosep,leftmargin=*]
  \item 定义目标期望或统计量、样本单位和依赖假设；
  \item 给出解析、枚举或退化参数 oracle；
  \item 推导误差率，声明样本量与统计容差来源；
  \item 记录 RNG、seed/substream、并行与软件环境；
  \item 比较方差缩减前后的同一目标，而非只比较方差；
  \item 对依赖、重尾、稀有事件和选择步骤运行负例；
  \item 报告结论能证明什么、不能外推什么。
\end{enumerate}

\section*{章末练习}
\begingroup\small
\subsection*{口述概念}
\begin{enumerate}[nosep,leftmargin=*]
  \item 为什么固定 seed 只能证明可重放，不能证明 Monte Carlo 正确？
  \item 区分模型误差、离散化偏差、抽样误差、数值误差和实现错误。
  \item Monte Carlo 与 bootstrap 都会重复抽样，它们分别在近似什么分布下的什么对象？
\end{enumerate}

\subsection*{笔试推导}
\begin{enumerate}[nosep,leftmargin=*]
  \item 从协方差双重和推导 $\operatorname{Var}(\hat\mu_N)=\sigma^2/N$，并说明相关样本时哪一步失效。
  \item 推导控制变量的 $\beta^*$ 与最小方差；对 $U^2$ 用 $U$ 作控制变量算出方差缩减倍数。
  \item 对样本 $(1,2,4)$ 用和的频数表推导 bootstrap 均值与方差。
\end{enumerate}

\subsection*{数值编程}
\begin{enumerate}[leftmargin=*]
  \item 比较普通、控制变量和重要性抽样的逐样本贡献与标准误。
  \item 改变控制变量中心，解释低方差为何仍可能错误。
  \item 穷举样本 $(1,2,4)$ 的全部有序重抽样均值。
\end{enumerate}

\subsection*{研究判断}
\begin{enumerate}[nosep,leftmargin=*]
  \item 有自相关的日收益逐日 bootstrap 后得到很窄的 Sharpe 区间，结论为何不可信？
  \item 重要性抽样产生巨大权重和极小的某次误差，应追加哪些证据？
  \item 因子从五百个候选中筛选后再 bootstrap 最终一个因子，区间遗漏了什么不确定性？
\end{enumerate}
\endgroup

\subsection*{基础衔接：独立计算}
重要性抽样采用 $q(x)=3x^2$ 时，积分 $\int_0^1x^2dx$ 的贡献和方差是多少？为什么现实中难以普遍实现？

\subsection*{分级提示}
\begingroup\small
\begin{enumerate}[nosep,leftmargin=*]
  \item \textbf{口述概念：}seed 管可重放性，标准误管条件于模型的抽样波动，独立基准管实现正确性；bootstrap 的随机性还条件于已经观察到的样本。
  \item \textbf{笔试推导：}均值方差先写协方差双重和；控制变量方差是关于 $\beta$ 的二次函数；bootstrap 小样本可对和的频数做完全枚举。
  \item \textbf{数值编程：}先写出本章公式中的目标量，再显示中间计算。改变参数前预测结果方向，运行后说明差异来自模型、抽样还是数值近似。
  \item \textbf{研究判断：}先识别真正独立的抽样单位，再选 block、cluster 或完整管线重抽样；重要性权重要查支撑、尾部和有效样本量。
\end{enumerate}
\endgroup


\section{章末产出：一张模拟审计卡}

为一个自己熟悉的模拟任务提交一页审计卡，而不是只提交 notebook。审计卡至少包含：目标 $\mu$ 的数学定义；抽样分布与独立单位；一个解析、枚举或退化参数 oracle；抽样、离散化与数值误差预算；RNG 与子流协议；方差缩减为何保持目标不变；一个会稳定失败的错误配置；结论的适用边界。

本章 $U^2$ 实验的最小版本可以写成：目标为 $\mathbb E[U^2]=1/3$；解析方差为 $4/45$；固定 seed 只用于复跑，四标准误用于统计验收；控制变量 $U$ 的已知期望为 $1/2$，均值与 VRF 必须同时通过；把中心改成 $0.6$ 应被拒绝；bootstrap 的 27 状态枚举只验证 IID 条件世界，不外推到相关收益。若审计卡无法说明其中任一项，程序输出就还不足以支撑研究结论。

\noindent\textbf{本章小结与 Capstone 连接}\quad
Monte Carlo 的 $N^{-1/2}$ 误差率来自方差推导；方差缩减必须同时保持目标期望；bootstrap 近似的是条件抽样分布，不能脱离依赖和正则性假设。上册 Capstone 将要求每项模拟结论附带解析/枚举 oracle、误差预算、随机流协议、方差缩减有效性与失败边界。
